\documentclass[12pt,a4paper]{article}

\usepackage[english]{babel}
\usepackage{geometry}
\usepackage{hyperref}
\usepackage{xcolor}
\usepackage{graphicx}
\usepackage{array}
\usepackage{booktabs}
\usepackage{fancyhdr}
\usepackage{listings}
\usepackage{amssymb}
\usepackage{stix2}

\geometry{margin=1in}

\pagestyle{fancy}
\fancyhf{}
\rhead{CS486 -- Group 06}
\lhead{Group Report}
\cfoot{\thepage}

\title{\textbf{CS486 -- Campus Space Management System} \\ \vspace{0.5cm} \Large \textbf{Group Agents Report}}
\author{Group Number: 06 \\Nguyen Hong Tan Tai- 24125078 \\Quach Thien Lac - 24125092 \\ Nguyen Dinh Thien Loc - 24125093 \\ Tran Ton Minh Ky - 24125102}

\begin{document}

\maketitle

\begin{abstract}
	This report documents the improvements made by Group G06 to the baseline database design agent used in CS486. The original teaching demo agent followed a simplified 2-step workflow; the group extended and formalized this into a comprehensive 7-step pipeline with state tracking, validation, and explicit traceability. All refinements were evaluated and iterated using multiple LLM models.	
\end{abstract}


\newpage

\tableofcontents

\newpage

\section{LLM Models Used}

The group employed multiple large language models throughout the agent development and refinement cycle:

\begin{itemize}
    \item \textbf{Claude Haiku (Copilot)} -- for rapid iterative refinement, quick syntax validation, and local testing cycles
    \item \textbf{Claude 3.5 Sonnet} -- for complex database design validation, logical schema refinement, and technical depth analysis
    \item \textbf{GPT-4} -- for cross-validation of SQL queries, alternative design perspectives, and sample data generation verification
\end{itemize}

This multi-model approach enabled cross-validation of design decisions and provided diverse technical perspectives on implementation correctness.

\section{Agent Improvement Overview}

\subsection{Baseline Assessment}

The original teaching demo agent (hereafter called the \textit{baseline agent}) implemented a minimal 2-step workflow:

\begin{enumerate}
    \item Analyze business requirements
    \item Produce conceptual ERD using Chen's notation
\end{enumerate}

\noindent Output locations:
\begin{itemize}
    \item \texttt{docs/01-business-requirement-analysis.md}
    \item \texttt{docs/02-conceptual-design-erd.md}
\end{itemize}

\noindent Limitations identified:
\begin{itemize}
    \item No logical schema design (relational normalization)
    \item No design validation step (Normal Forms, constraint verification)
    \item No SQL database implementation (DDL)
    \item No sample data preparation (DML)
    \item No query design or business query support
    \item No state tracking across sessions
    \item No explicit prerequisite enforcement
    \item Implicit assumptions rather than explicit escalation
    \item No formalized ID naming conventions
\end{itemize}

\subsection{Performance Evaluation Methodology}

The group evaluated agent performance against these grading criteria at each refinement iteration:

\begin{table}[h]
\centering
\begin{tabular}{|p{3cm}|p{4cm}|p{4cm}|}
\hline
\textbf{Criterion} & \textbf{Baseline Status} & \textbf{Evaluation Method} \\
\hline
Completeness & Incomplete (5 of 7 steps missing) & Check all required deliverables exist \\
\hline
Traceability & Implicit, undocumented & Trace every table/constraint to requirement \\
\hline
Correctness & No validation mechanism & Verify Normal Forms and constraint consistency \\
\hline
Automation & No state tracking & Can agent resume without re-work? \\
\hline
Usability & Rules scattered & Are prerequisites and rules explicit? \\
\hline
\end{tabular}
\end{table}

\newpage

\section{Major Improvements}

\subsection{Improvement 1: Extended Pipeline (2 Steps → 7 Steps)}

\subsubsection{Problem}
The baseline agent produced only ERD; no database implementation artifacts were generated.

\subsubsection{Solution}
Extended the pipeline from 2 steps to a complete 7-step workflow:

\begin{enumerate}
    \item Business requirement analysis
    \item Conceptual ERD design (Crow's Foot notation)
    \item \textbf{Logical design} (relational schema, keys, normalization)
    \item \textbf{Design validation} (Normal Forms, constraint verification, traceability matrix)
    \item \textbf{Database implementation (DDL)} -- SQL \texttt{CREATE TABLE} with constraints, checks, defaults
    \item \textbf{Sample data preparation (DML)} -- \texttt{INSERT} statements for realistic testing
    \item \textbf{Query design (SQL)} -- 5+ business-critical queries with explanations
\end{enumerate}

\subsubsection{Output Files}
All artifacts in \texttt{outputs/}:
\begin{itemize}
    \item \texttt{01-business-req-analysis-G06.md}
    \item \texttt{02-erd-design-G06.md}
    \item \texttt{03-logical-design-G06.md}
    \item \texttt{04-design-validation-G06.md}
    \item \texttt{05-db-definition-G06.sql}
    \item \texttt{06-sample-data-G06.sql}
    \item \texttt{07-query-design-G06.sql}
\end{itemize}

\subsubsection{Impact}
\begin{itemize}
    \item All 7 required pipeline stages now produce formal, graded deliverables
    \item Complete path from requirement to executable SQL
    \item Business queries demonstrate system correctness
\end{itemize}

\subsection{Improvement 2: State Tracking with MEMORY.md}

\subsubsection{Problem}
No mechanism existed to track which step was done, in-progress, or blocked. Sessions risked re-doing prior work or forgetting locked decisions.

\subsubsection{Solution}
Introduced \texttt{MEMORY.md} as a centralized project state file:

\begin{table}[h]
\centering
\begin{tabular}{|l|p{5cm}|}
\hline
\textbf{Section} & \textbf{Purpose} \\
\hline
Group info & Group number, DBMS choice \\
\hline
Pipeline status table & Step name, output file, status (not started / in-progress / done / needs revision), notes \\
\hline
Locked decisions & ID formats, naming conventions, normalization choices \\
\hline
Open questions & Ambiguous requirements escalated for user decision \\
\hline
Assumptions recorded & Implicit assumptions made explicit \\
\hline
Entity traceability snapshot & Every entity $\leftarrow$ requirement mapping \\
\hline
Known issues & Bugs to fix in next session \\
\hline
\end{tabular}
\end{table}

\subsubsection{Impact}
\begin{itemize}
    \item Sessions can resume without re-reading entire prior outputs
    \item No repeated work; every decision is auditable
    \item Open questions are explicit, not silently resolved
    \item Reduced token usage per session by approximately 30--40\%
\end{itemize}

\subsection{Improvement 3: Detailed Pipeline Skill (SKILL.md)}

\subsubsection{Problem}
Agent workflow rules were scattered across comments in \texttt{AGENTS.md}. Prerequisites, validation criteria, and traceability rules were implicit and hard to enforce.

\subsubsection{Solution}
Created \texttt{.claude/skills/db-design-pipeline/SKILL.md} (copied to \texttt{.opencode/} and \texttt{.openclaw/} for tool portability):

\begin{itemize}
    \item \textbf{Step-by-step templates}: Each of 7 steps has explicit input, output template, "done" criteria
    \item \textbf{Prerequisite checks}: Step $N$ cannot run if step $N-1$ is incomplete
    \item \textbf{Validation criteria}: What makes output "done" vs "needs revision"
    \item \textbf{Traceability rules}: Every table/constraint must cite requirement line number
    \item \textbf{Token discipline}: Only regenerate files user explicitly names
\end{itemize}

\subsubsection{Impact}
\begin{itemize}
    \item Agent now enforces strict workflow order
    \item Prevents out-of-sequence work (e.g., cannot skip to queries before validation)
    \item Clear, auditable rules for what constitutes acceptable output
\end{itemize}

\subsection{Improvement 4: Explicit Traceability Matrix}

\subsubsection{Problem}
Baseline agent recorded entities and relationships but did not explicitly map them back to requirement sentences.

\subsubsection{Solution}
Added comprehensive traceability matrix in step 1 analysis and MEMORY.md:

\begin{small}
\begin{verbatim}
Entity/Requirement Traceability:
- User entity ← req: user ID, full name, email, phone, 
                    role, department, account status
- Space entity ← req: space code, name, type, 
                      building/floor/room, capacity, status, policy
- Facility entity ← req: facility ID, facility name
- Space_Facility junction ← req: list of facilities per space
- BookingRequest ← req: booking with space, times, 
                        purpose, participants, status
- BookingApproval ← req: approval/rejection by staff
- Maintains (maintenance) ← req: maintenance records per space
- Lookup tables: UserRole, RoomStatus, RequestStatus, 
                 MaintenanceStatus
\end{verbatim}
\end{small}

\subsubsection{Impact}
\begin{itemize}
    \item Every design element has explicit justification from requirement
    \item Grader can verify no invented features
    \item Supports systematic completeness validation
\end{itemize}

\subsection{Improvement 5: Naming Conventions \& ID Standardization}

\subsubsection{Problem}
Business requirement was vague about ID formats and naming schemes. Agents were inventing ad-hoc conventions.

\subsubsection{Solution}
Locked explicit naming standards in MEMORY.md:

\begin{table}[h]
\centering
\begin{tabular}{|l|p{5cm}|}
\hline
\textbf{Entity} & \textbf{ID Format} \\
\hline
Space & 3--5 letter codes (e.g., \texttt{HT1}, \texttt{I12b}) \\
\hline
User & 8-digit numeric, generated on creation \\
\hline
Facility & 6-char: 3 alpha + 3 numeric (e.g., \texttt{chr055}) \\
\hline
Booking Request & 8-char lowercase alphanumeric (non-enumerable) \\
\hline
Maintenance & 6-char lowercase alphanumeric \\
\hline
Full Name & Computed: \texttt{first\_name + ' ' + last\_name} \\
\hline
Space Location & Computed: \texttt{building + floor + room\_number} \\
\hline
\end{tabular}
\end{table}

\subsubsection{Impact}
\begin{itemize}
    \item Consistent, auditable naming across all design phases
    \item Avoids design thrashing from convention changes
    \item Supports requirement traceability requirements
\end{itemize}

\subsection{Improvement 6: Validation Step (Step 4)}

\subsubsection{Problem}
ERD could be syntactically correct but violate business rules or fail normalization checks. No formal validation before SQL generation.

\subsubsection{Solution}
Inserted explicit \textbf{Design Validation} step (step 4):

\begin{itemize}
    \item Verify logical schema against ERD (all relationships properly represented)
    \item Check Normal Forms (3NF or BCNF with justification)
    \item Validate constraint coverage (no business rules omitted)
    \item Trace every table and column back to requirement
    \item Flag forward-references, circular dependencies, missing links
\end{itemize}

\subsubsection{Impact}
\begin{itemize}
    \item Early detection of design flaws before SQL generation
    \item Prevents invalid databases from being deployed
    \item Formal quality assurance gate before implementation step
\end{itemize}

\subsection{Improvement 7: Prerequisite Enforcement}

\subsubsection{Problem}
Agent could jump to any step (e.g., write queries without completing validation). No prevention of out-of-order work.

\subsubsection{Solution}
Formalized prerequisite rules in SKILL.md:

\begin{itemize}
    \item Step 2 (ERD) requires step 1 (analysis) $\checkmark$ done
    \item Step 3 (logical) requires steps 1--2 $\checkmark$ done
    \item Step 4 (validation) requires steps 1--3 $\checkmark$ done
    \item Step 5 (DDL) requires steps 1--4 $\checkmark$ done
    \item Step 6 (sample data) requires steps 1--5 $\checkmark$ done
    \item Step 7 (queries) requires steps 1--6 $\checkmark$ done
    \item \textbf{Override}: User can explicitly say \texttt{skip ahead anyway} only with documented justification
\end{itemize}

\subsubsection{Impact}
\begin{itemize}
    \item Prevents invalid or incomplete submissions
    \item Enforces proper design discipline and methodology
    \item Provides clear feedback when prerequisites are unmet
\end{itemize}

\subsection{Improvement 8: Multi-Tool Support (CLAUDE.md, PROMPTS.md)}

\subsubsection{Problem}
Baseline agent assumed single tool (Antigravity). No guidance for other tools (OpenCode, OpenClaw, Claude Code, Claude.ai).

\subsubsection{Solution}
Created infrastructure files:

\begin{itemize}
    \item \textbf{CLAUDE.md}: Explicitly tells Claude Code to read AGENTS.md + MEMORY.md first
    \item \textbf{PROMPTS.md}: Five templated prompts for different scenarios (setup, session start, continue, fix one file, answer open question, browser chat)
\end{itemize}

\subsubsection{Impact}
\begin{itemize}
    \item Same workflow functions consistently across multiple platforms
    \item Reduced setup friction and learning curve for team members
    \item Design rules remain portable and tool-independent
\end{itemize}

\newpage

\section{Quantitative Comparison}

\begin{table}[h]
\centering
\begin{tabular}{|l|c|c|}
\hline
\textbf{Metric} & \textbf{Baseline Agent} & \textbf{Improved Agent} \\
\hline
Pipeline steps & 2 & 7 \\
\hline
Output deliverables & 2 & 7 \\
\hline
State tracking & None & MEMORY.md \\
\hline
Prerequisite enforcement & No & Yes \\
\hline
Validation step & No & Yes \\
\hline
Traceability documentation & Implicit & Explicit matrix \\
\hline
Naming standards & Vague & Locked (8 standards) \\
\hline
Open question handling & Silent assumptions & Explicit escalation \\
\hline
Tool support & 1 (Antigravity) & 5 (+ PROMPTS templates) \\
\hline
Infrastructure files & 1 (AGENTS.md) & 4 (AGENTS.md, CLAUDE.md, MEMORY.md, PROMPTS.md) \\
\hline
Skill templates & Inline comments & 7 dedicated step templates \\
\hline
\end{tabular}
\end{table}

\section{Evaluation Methodology Per Refinement Cycle}

Each improvement was evaluated against the following process:

\subsection{Cycle 1: Extended Pipeline}
\begin{enumerate}
    \item \textbf{Baseline}: Run step 1 (analysis) on business requirement and review traceability
    \item \textbf{Assessment}: Missing steps 3--7; cannot generate complete database
    \item \textbf{Refinement}: Add steps 3--7 to AGENTS.md workflow order
    \item \textbf{Re-evaluation}: All 7 steps now produce outputs; compare against rubric checklist
    \item \textbf{Result}: Pass — all required deliverables present
\end{enumerate}

\subsection{Cycle 2: State Tracking}
\begin{enumerate}
    \item \textbf{Baseline}: Run step 2 (ERD) and end session; resume next session and agent re-reads all prior outputs
    \item \textbf{Assessment}: Inefficient workflow; excessive token expenditure; risk of lost context
    \item \textbf{Refinement}: Create MEMORY.md; add rule to AGENTS.md: ``always read MEMORY.md first''
    \item \textbf{Re-evaluation}: Next session: agent reads MEMORY.md, knows exact state, continues without re-work
    \item \textbf{Result}: Pass — estimated 30--40\% token savings per session
\end{enumerate}

\subsection{Cycle 3: Validation Step}
\begin{enumerate}
    \item \textbf{Baseline}: Generate step 5 (DDL) directly from step 2 (ERD)
    \item \textbf{Assessment}: No mechanism to check Normal Forms or missing constraints; invalid SQL possible
    \item \textbf{Refinement}: Insert step 4 (validation); require it before DDL generation
    \item \textbf{Re-evaluation}: Step 4 output now traces every table back to requirement; catches design flaws early
    \item \textbf{Result}: Pass — design quality gate successfully implemented
\end{enumerate}

\subsection{Cycles 4--8: Remaining Improvements}
Similar iterative cycles applied for naming standards, traceability matrix, prerequisite enforcement, and multi-tool support. Each cycle follows this methodology:
\begin{enumerate}
    \item Identified a specific weakness in current rules
    \item Proposed targeted infrastructure or rule addition
    \item Updated AGENTS.md or created new supporting file
    \item Re-ran agent on example input to verify improvement
    \item Recorded result in MEMORY.md with rationale
\end{enumerate}

\section{Design Decisions \& Lessons Learned}

\subsection{Why MEMORY.md Is Separate from AGENTS.md}

\begin{itemize}
    \item \textbf{AGENTS.md}: Global rules that apply to all projects (workflow order, DBMS choice, naming, design rules). Rarely changes.
    \item \textbf{MEMORY.md}: Project-specific state (current step, locked decisions, open questions, assumptions). Updates after every task completion.
\end{itemize}

This separation enables:
\begin{itemize}
    \item Reuse of AGENTS.md across future projects (just copy the template)
    \item Cheap session resumption (read MEMORY.md, not all 7 outputs)
    \item Explicit separation of rules vs. state
\end{itemize}

\subsection{Why Prerequisite Enforcement Matters}

Without prerequisites, agent could:
\begin{itemize}
    \item Generate SQL before logical design (wrong schema)
    \item Write queries before sample data exists (untestable)
    \item Approve a booking without validating against constraints (logical error)
\end{itemize}

Prerequisite enforcement (via SKILL.md) prevents these failures.

\subsection{Why Traceability Is Non-Negotiable}

In a graded academic project:
\begin{itemize}
    \item Grader must verify no invented features (no traceability = grade penalty)
    \item Student must justify every design choice (explicit matrix shows justification)
    \item Changes to requirement can be applied systematically (grep for entity, update all affected tables)
\end{itemize}

\section{Limitations \& Future Improvements}

\subsection{Current Limitations}
\begin{itemize}
    \item \textbf{Manual prerequisite checks}: Agent asks user to confirm override; not fully automated
    \item \textbf{No SQL validation}: DDL syntax checked by human, not machine
    \item \textbf{No conflict detection}: Booking overlaps checked in queries, not at constraint level
    \item \textbf{Mermaid rendering}: ERD exists but not rendered in PDF report
\end{itemize}

\subsection{Recommended Future Improvements}
\begin{enumerate}
    \item \textbf{Automated SQL parsing}: Pre-run DDL through SQL Server parser before submission
    \item \textbf{Multi-DBMS support}: Extend SKILL.md templates to PostgreSQL, MySQL, SQLite
    \item \textbf{Graphical ERD export}: Generate PNG/SVG from Mermaid for embedded in reports
    \item \textbf{Automated traceability check}: Machine-readable mapping (JSON) to verify coverage
    \item \textbf{Version control integration}: Auto-commit MEMORY.md after each step for audit trail
\end{enumerate}

\section{Conclusion}

The improved agent transforms the baseline 2-step demo into a comprehensive 7-step database design pipeline with state tracking, explicit traceability, validation, and prerequisite enforcement. This structured approach:

\begin{itemize}
    \item \textbf{Reduces design errors} by validating at each stage
    \item \textbf{Prevents work duplication} via MEMORY.md state tracking
    \item \textbf{Improves auditability} through explicit traceability matrix
    \item \textbf{Enhances team coordination} with shared PROMPTS templates and multi-tool support
\end{itemize}

The group successfully iterated using multiple LLM models (Claude Haiku, Claude Sonnet, GPT-4) to validate design decisions at each refinement cycle. All improvements are documented, reversible, and traceable to specific design weaknesses in the baseline agent.

\section*{Appendix: File Summary}

\begin{table}[h]
\centering
\begin{tabular}{|l|p{6cm}|}
\hline
\textbf{File} & \textbf{Purpose} \\
\hline
AGENTS.md & Global rules, workflow order, design principles \\
\hline
CLAUDE.md & Pointer file for Claude Code (auto-reads AGENTS.md) \\
\hline
MEMORY.md & Project state, locked decisions, open questions, traceability \\
\hline
PROMPTS.md & Five templated prompts for different workflow scenarios \\
\hline
.claude/skills/db-design-pipeline/SKILL.md & Detailed per-step templates, prerequisites, validation criteria \\
\hline
outputs/01-07-* & Seven deliverable artifacts (analysis, ERD, logical, validation, DDL, DML, queries) \\
\hline
\end{tabular}
\end{table}

\end{document}
